% P46 Figure: p(L,C,R) vs V_Phi(Phi) — vertical stack
\begin{figure}[htbp]
\centering
\begin{tikzpicture}[
  scale=0.95,
  every node/.style={font=\small},
  colbox/.style={draw,rounded corners=4pt,align=left,inner sep=8pt,text width=6.2cm}
]
  \node[colbox,fill=orange!8,draw=orange!55!black] (pbox) at (0,2.35) {
    \textbf{$p(L,C,R)$} --- update operator\\[4pt]
    \scriptsize
    $\mathrm{GF}(7)$ neighbourhood update\\
    gravity: $\nabla^2\Phi \propto p(w_x,w_y,w_z)$\\
    chiral; cubic term $-LCR$\\
    Rule~110 on binary projection
  };
  \node[colbox,fill=green!8,draw=green!45!black] (vbox) at (0,-2.35) {
    \textbf{$V_{\ZS}(\Phi)$} --- potential energy\\[4pt]
    \scriptsize
    $\dfrac{m^2}{49}(1-\cos 7\Phi)$\\
    seven vacua $\Phi_k=2\pi k/7$\\
    BPS kinks and domain walls\\
    smooth continuum Lagrangian
  };
  \draw[ugparr] (pbox.south) -- (vbox.north)
    node[midway,right=6pt,font=\scriptsize,align=left]
    {same $F_{21}$ symmetry\\[-2pt]different mathematical role};
\end{tikzpicture}
\caption{The framework contains two distinct $\mathbb{Z}_7$-periodic objects:
  the discrete update polynomial (dynamics and coupling source) and the continuum
  potential (vacuum structure and solitons). Neither replaces the other.}
\label{fig:p_vs_Vphi}
\end{figure}
